\chapter{สถาปัตยกรรมโมเดลและการเรียนรู้เพิ่มเติมแบบทันที}

\section{โครงสร้างโมเดลการเรียนรู้เชิงลึกแบบสองหัว (Dual-Head Architecture)}
ระบบ DurianLeaf AI ใช้โครงข่ายประสาทเทียมแบบคอนโวลูชันชนิดประสิทธิภาพสูง (MobileNetV3 Backbone) รับข้อมูลภาพขนาด $224 \times 224 \times 3$ พิกเซล และแยกผลลัพธ์ออกเป็นสองหัวการทำนายแบบขนาน ดังแสดงในรูปที่ \ref{fig:dual_head_architecture}

\begin{figure}[h]
\centering
\begin{tikzpicture}[
  box/.style={draw=rbruDarkGreen, fill=white, line width=1pt, rounded corners=4pt, align=center, font=\small},
  head/.style={draw=rbruForestGreen, fill=rbruForestGreen!10, line width=1.2pt, rounded corners=6pt, align=center, font=\small\bfseries},
  arrow/.style={->, >=stealth, line width=1.2pt, color=rbruDarkGreen}
]
  \node[box, fill=black!5] (input) at (0, 0) {ภาพใบไม้ทุเรียนนำเข้า\\$224 \times 224 \times 3$};
  \node[box, fill=rbruDarkGreen!15] (backbone) at (0, -1.5) {โครงข่ายหลัก MobileNetV3\\Depthwise Separable Convolutions};
  \node[box, fill=rbruGold!20] (embed) at (0, -3.0) {ชั้นเวกเตอร์เอกลักษณ์ (Embedding Layer)\\$128\text{-Dimensional Latent Vector}$};

  \node[head] (head1) at (-3.5, -4.8) {หัวทำนายที่ 1 (Pathology)\\Softmax Classifier\\6 คลาสโรคและศัตรูพืช};
  \node[head] (head2) at (3.5, -4.8) {หัวทำนายที่ 2 (Nutrients)\\Linear Regressor\\9 ตัวแปร NPK และ SPAD};

  \draw[arrow] (input) -- (backbone);
  \draw[arrow] (backbone) -- (embed);
  \draw[arrow] (embed) -| (head1);
  \draw[arrow] (embed) -| (head2);
\end{tikzpicture}
\caption{แผนผังสถาปัตยกรรมโมเดล Dual-Head MobileNet และชั้นเวกเตอร์เอกลักษณ์}
\label{fig:dual_head_architecture}
\end{figure}

\section{กลไกการเรียนรู้เพิ่มเติมแบบทันที (On-Device Continual Learning)}
หนึ่งในข้อจำกัดของโมเดล Deep Learning ทั่วไปคือเมื่อนำไปติดตั้งบนอุปกรณ์เคลื่อนที่แล้ว จะไม่สามารถเรียนรู้โรคใหม่หรืออาการเฉพาะถิ่นได้ ระบบ DurianLeaf AI แก้ไขปัญหานี้ด้วยกลไก Few-Shot Metric Learning

เมื่อผู้ใช้งานบันทึกภาพตัวอย่างอาการใหม่ $3 - 5$ ภาพ โมเดลจะส่งภาพผ่านโครงข่ายหลักเพื่อสกัดเวกเตอร์เอกลักษณ์ขนาด 128 มิติ ($\mathbf{u} \in \mathbb{R}^{128}$) และทำการปรับมาตรฐาน (L2 Normalization) เก็บไว้ในหน่วยความจำเครื่อง

เมื่อกล้องตรวจพบภาพใบไม้ใหม่ ระบบจะคำนวณความคล้ายคลึงโคไซน์ (Cosine Similarity) ระหว่างเวกเตอร์ภาพใหม่ $\mathbf{v}$ กับเวกเตอร์อ้างอิง $\mathbf{u}$
\begin{equation}
  \text{Cosine Similarity}(\mathbf{u}, \mathbf{v}) = \frac{\mathbf{u} \cdot \mathbf{v}}{\|\mathbf{u}\|_2 \|\mathbf{v}\|_2}
\end{equation}
หากค่าความคล้ายคลึงสูงกว่าเกณฑ์ที่กำหนด ($>0.85$) ระบบจะจำแนกเป็นอาการเฉพาะที่ผู้ใช้สอนได้ทันทีโดยไม่ต้องทำการฝึกสอนโครงข่ายใหม่ทั้งหมด

\section{วงจรการปรับปรุงโมเดลรอบใหญ่ (Active Learning Feedback Loop)}
สำหรับเคสที่มีความไม่แน่นอนสูง ระบบจะจัดเก็บข้อมูลลงในคิวบันทึก (Local Staging Queue) เพื่อส่งต่อไปยังชุดสคริปต์ Python ในเครื่องแม่ข่าย ทำการ Fine-Tuning และส่งออกเป็นโมเดลเวอร์ชันใหม่ผ่านระบบ Over-The-Air (OTA) ได้อย่างต่อเนื่อง
